\clearpage
\begin{landscape}
\begingroup
\fontsize{10}{13}\selectfont
\fancyhead[L]{\small POSTHUMOUS PAPERS. PERTURBATIONS OF PALLAS}
\renewcommand{\headrulewidth}{0pt}
\setlength{\lineskip}{4pt}\setlength{\lineskiplimit}{4pt}
\setlength{\arraycolsep}{4pt}
\makeatletter\tagsleft@true\makeatother
\newcommand{\VFivePair}{\mathord{\vcenter{\offinterlineskip\hbox{\rule{.12em}{.12em}}\kern.30em\hbox{\rule{.12em}{.12em}}}}}
\newcommand{\VFivePoint}{\mathord{.}\mkern5mu}
\newcommand{\VFiveTenPoints}{\VFivePoint\VFivePoint\VFivePoint\VFivePoint\VFivePoint\VFivePoint\VFivePoint\VFivePoint\VFivePoint\VFivePoint}
\pdfbookmark[2]{General perturbations, first calculation, [5]}{pallas-V-fragment5}
\SourceLeaf{511}
\begin{center}[5.]\end{center}
\par\noindent\hspace*{\parindent}[If a function of two arguments of the form
\[
\begin{aligned}
T&=\gamma_{0,0}+\gamma_{1,0}\cos x+\gamma_{2,0}\cos2x+\cdots\\
 &\quad+\delta_{1,0}\sin x+\delta_{2,0}\sin2x+\cdots\\[4pt]
 &\quad+\gamma_{0,1}\cos y+\gamma_{1,1}\cos(x+y)+\gamma_{2,1}\cos(2x+y)+\cdots\\
 &\quad+\delta_{0,1}\sin y+\delta_{1,1}\sin(x+y)+\delta_{2,1}\sin(2x+y)+\cdots\\[4pt]
 &\qquad\qquad\VFivePair+\epsilon_{1,1}\cos(x-y)+\epsilon_{2,1}\cos(2x-y)+\cdots\\
 &\qquad\qquad+\zeta_{1,1}\sin(x-y)+\zeta_{2,1}\sin(2x-y)+\cdots\\[4pt]
 &\quad+\gamma_{0,2}\cos2y+\gamma_{1,2}\cos(x+2y)+\gamma_{2,2}\cos(2x+2y)+\cdots\\
 &\quad+\delta_{0,2}\sin2y+\delta_{1,2}\sin(x+2y)+\delta_{2,2}\sin(2x+2y)+\cdots\\[4pt]
 &\qquad\qquad+\epsilon_{1,2}\cos(x-2y)+\epsilon_{2,2}\cos(2x-2y)+\cdots\\
 &\qquad\qquad+\zeta_{1,2}\sin(x-2y)+\zeta_{2,2}\sin(2x-2y)+\cdots\\[4pt]
 &\quad+\VFiveTenPoints\VFiveTenPoints\VFiveTenPoints\VFiveTenPoints\VFivePoint\VFivePoint\VFivePoint\VFivePoint\VFivePoint\VFivePoint\VFivePoint
\end{aligned}\tag*{1)}
\]
is to be expanded, the procedure explained in the preceding articles can also be applied. The function \(T\) can be written in the following form:
\[
\begin{aligned}
T&=\alpha_0+\alpha_1\cos x+\alpha_2\cos2x+\cdots\\
 &\quad+\beta_1\sin x+\beta_2\sin2x+\cdots
\end{aligned}\tag*{2)}
\]
where the \(\alpha_0,\ \alpha_1,\ \alpha_2\ldots\beta_1,\ \beta_2\ldots\) are functions of \(y\) alone, which can in turn be represented in the form
\[
\begin{aligned}
\alpha_0&=p_{0,0}+p_{0,1}\cos y+p_{0,2}\cos2y+\cdots\\
 &\quad+q_{0,1}\sin y+q_{0,2}\sin2y+\cdots\\[4pt]
\alpha_1&=p_{1,0}+p_{1,1}\cos y+p_{1,2}\cos2y+\cdots\\
 &\quad+q_{1,1}\sin y+\VFivePair q_{1,2}\sin2y+\cdots\\
 &\text{etc.}\\[4pt]
\beta_1&=r_{1,0}+r_{1,1}\cos y+r_{1,2}\cos2y+\cdots\\
 &\quad+s_{1,1}\sin y+s_{1,2}\sin2y+\cdots\\
 &\text{etc.}
\end{aligned}\tag*{3)}
\]
with the \(p_{0,0},\ p_{0,1},\ldots q_{0,1},\ q_{0,2},\ldots r_{1,0},\ r_{1,1},\ldots s_{1,1},\ s_{1,2},\ldots\) being constant coefficients.

\par\noindent\hspace*{\parindent}Suppose now that for the values \(x=0,\ a,\ 2a,\ldots(\mu-1)a\) and \(y=0,\ a',\ 2a',\ldots(\nu-1)a'\) the corresponding values of the function \(T\) are given, with \(T=A_{0,0}\) for \(x=0,\ y=0\), \(T=A_{1,0}\) for \(x=a,\ y=0\), \(T=A_{2,0}\) for \(x=2a,\ y=0\), \(\ldots T=A_{0,1}\) for \(x=0,\ y=a'\), and \(\ldots T=A_{\mu-1,\nu-1}\) for \(x=(\mu-1)a,\ y=(\nu-1)a'\). Thus in all \(\mu\nu\) values of \(T\) are given. As in the preceding articles, we shall assume here that \(\mu a=\nu a'=360^0\) and that both \(\mu\) and \(\nu\) are even. Let \(T_0\) denote the value taken by \(T\) when \(y=0\), \(T_1\) the value of \(T\) for \(y=a'\), \(T_2\) that for \(y=2a'\), \(\ldots\), and \(T_{\nu-1}\) that for \(y=(\nu-1)a'\). Thus \(T_0,\ T_1\) etc. are functions of \(x\) alone, which can in turn be
\SourceLeaf{512}
\[
\begin{aligned}
T_0&=\alpha_{0,0}+\alpha_{1,0}\cos x+\alpha_{2,0}\cos2x+\cdots\\
 &\quad+\beta_{1,0}\sin x+\beta_{2,0}\sin2x+\cdots\\[4pt]
T_1&=\alpha_{0,1}+\alpha_{1,1}\cos x+\alpha_{2,1}\cos2x+\cdots\\
 &\quad+\beta_{1,1}\sin x+\beta_{2,1}\sin2x+\cdots\\[4pt]
 &\VFiveTenPoints\VFiveTenPoints\VFiveTenPoints\VFivePoint\VFivePoint\\[4pt]
T_{\nu-1}&=\alpha_{0,\nu-1}+\alpha_{1,\nu-1}\cos x+\alpha_{2,\nu-1}\cos2x+\cdots\\
 &\quad+\beta_{1,\nu-1}\sin x+\beta_{2,\nu-1}\sin2x+\cdots
\end{aligned}\tag*{4)}
\]
written in this form.

\par\noindent\hspace*{\parindent}Clearly \(A_{0,0},\ A_{1,0},\ A_{2,0}\ldots A_{\mu-1,0}\) are, in order, the values taken by \(T_0\) for \(x=0,\ a,\ 2a,\ldots(\mu-1)a\). From these, following precisely the prescriptions of Articles [2]---[4], one can determine the coefficients \(\alpha_{0,0},\ \alpha_{1,0},\ldots\beta_{1,0},\ \beta_{2,0}\ldots\). Thus, e.~g.,
\[
\begin{aligned}
\mu\alpha_{0,0}&=A_{0,0}+A_{1,0}+A_{2,0}+\cdots A_{\mu-1,0}\\[4pt]
\frac\mu2\alpha_{1,0}&=A_{0,0}+A_{1,0}\cos a+A_{2,0}\cos2a+\cdots+A_{\mu-1,0}\cos(\mu-1)a\\[4pt]
\frac\mu2\beta_{1,0}&=A_{0,0}+A_{1,0}\sin a+A_{2,0}\sin2a+\cdots+A_{\mu-1,0}\sin(\mu-1)a\\
 &\text{etc.}
\end{aligned}
\]

\par\noindent\hspace*{\parindent}Likewise \(A_{0,1},\ A_{1,1},\ A_{2,1}\ldots A_{\mu-1,1}\) are, in order, the values of \(T_1\) for \(x=0,\ a,\ 2a,\ldots(\mu-1)a\), from which the coefficients \(\alpha_{0,1},\ \alpha_{1,1},\ \alpha_{2,1}\ldots\beta_{1,1},\ \beta_{2,1}\ldots\) can be determined in the same way. The same procedure is to be applied to \(T_2,\ T_3,\ldots T_{\nu-1}\).

\par\noindent\hspace*{\parindent}Once these calculations have been carried out, expressions 4 for the functions \(T_0,\ T_1,\ldots T_{\nu-1}\) are completely determined.

\par\noindent\hspace*{\parindent}Furthermore, the now known coefficients \(\alpha_{0,0},\ \alpha_{0,1},\ \alpha_{0,2}\ldots\alpha_{0,\nu-1}\) are, in order, the values taken by \(\alpha_0\) for \(y=0,\ a',\ 2a'\ldots(\nu-1)a'\), and \(\alpha_{1,0},\ \alpha_{1,1},\ \alpha_{1,2},\ldots\alpha_{1,\nu-1}\) are the values of \(\alpha_1\) for \(y=0,\ a',\ 2a'\ldots(\nu-1)a'\), etc. Likewise \(\beta_{1,0},\ \beta_{1,1},\ \beta_{1,2},\ldots\beta_{1,\nu-1}\) are the values of \(\beta_1\) for the stated values of \(y\). Thus one can determine
\[
\begin{array}{@{}ll@{\quad}l@{\quad}l@{\quad}l@{}}
\text{the}&\text{coefficients}&p_{0,0},\ p_{0,1},\ p_{0,2},\ldots q_{0,1},\ q_{0,2}\ldots&\text{from}&\alpha_{0,0},\ \alpha_{0,1},\ \alpha_{0,2}\ldots\\[4pt]
\text{»}&\text{»}&p_{1,0},\ p_{1,1},\ p_{1,2},\ldots q_{1,1},\ q_{1,2},\ldots&\text{from}&\alpha_{1,0},\ \alpha_{1,1},\ \alpha_{1,2}\ldots\\[4pt]
\multicolumn{5}{@{}c@{}}{\VFiveTenPoints\VFiveTenPoints\VFiveTenPoints\VFiveTenPoints\VFiveTenPoints\VFivePoint\VFivePoint\VFivePoint\VFivePoint\VFivePoint}\\[4pt]
\text{»}&\text{»}&r_{1\,0},\ r_{1,1},\ r_{1,2},\ldots s_{1,1},\ s_{1,2},\ldots&\text{from}&\beta_{1,0},\ \beta_{1,1},\ \beta_{1,2}\ldots\\
\multicolumn{5}{c}{\text{etc.}}
\end{array}
\]
by again following the preceding articles.

\par\noindent\hspace*{\parindent}The expansions 3 for the \(\alpha_0,\ \alpha_1,\ \alpha_2,\ldots\beta_1,\ \beta_2,\ldots\) are now also completely known. By substituting these values in series 2 for \(T\), one obtains the complete expansion of the function \(T\) in both variables. The products of trigonometric functions that arise are converted into single sines and cosines of the sums and differences to obtain \(T\) in form 1.

\par\noindent\hspace*{\parindent}We may also give the formulas for calculating \(\gamma_{0,0},\allowbreak\ \gamma_{1,0},\allowbreak\ \gamma_{2,0},\allowbreak\ldots\delta_{1,0},\allowbreak\ \delta_{2,0},\allowbreak\ldots\epsilon_{1,1},\allowbreak\ \epsilon_{2,1},\allowbreak\ldots\zeta_{1,1},\allowbreak\ \zeta_{2,1},\allowbreak\ldots\) from \(p_{0,0},\allowbreak\ p_{0,1},\allowbreak\ldots q_{0,1},\allowbreak\ q_{0,2},\allowbreak\ldots r_{1,0},\allowbreak\ r_{1,1},\allowbreak\ldots s_{1,1},\allowbreak\ s_{1,2},\allowbreak\ldots\), although the calculation can be performed just as conveniently in a purely mechanical way:
\[
\begin{aligned}
\gamma_{0,0}&=p_{0,0},&\gamma_{0,1}&=p_{0,1},&\gamma_{0,2}&=p_{0,2}\quad\text{etc.}\\
&&\delta_{0,1}&=q_{0,1},&\delta_{0,2}&=q_{0,2}\quad\text{etc.}
\end{aligned}
\]
\SourceLeaf{513}
\[
\begin{aligned}
\gamma_{1,0}&=p_{1,0},&\gamma_{1,1}&=\tfrac12(p_{1,1}-s_{1,1}),&\gamma_{1,2}&=\tfrac12(p_{1,2}-s_{1,2})\quad\text{etc.}\\
\delta_{1,0}&=r_{1,0},&\delta_{1,1}&=\tfrac12(r_{1,1}+q_{1,1}),&\delta_{1,2}&=\tfrac12(r_{1,2}+q_{1,2})\quad\text{etc.}\\[8pt]
&&\epsilon_{1,1}&=\tfrac12(p_{1,1}+s_{1,1}),&\epsilon_{1,2}&=\tfrac12(p_{1,2}+s_{1,2})\quad\text{etc.}\\
&&\zeta_{1,1}&=\tfrac12(r_{1,1}-q_{1,1}),&\zeta_{1,2}&=\tfrac12(r_{1,2}-q_{1,2})\quad\text{etc.}\\[8pt]
\gamma_{2,0}&=p_{2,0},&\gamma_{2,1}&=\tfrac12(p_{2,1}-s_{2,1}),&\gamma_{2,2}&=\tfrac12(p_{2,2}-s_{2,2})\quad\text{etc.}\\
\delta_{2,0}&=r_{2,0},&\delta_{2,1}&=\tfrac12(r_{2,1}+q_{2\,1}),&\delta_{2,2}&=\tfrac12(r_{2,2}+q_{2,2})\quad\text{etc.}\\[8pt]
&&\epsilon_{2,1}&=\tfrac12(p_{2,1}+s_{2,1}),&\epsilon_{2,2}&=\tfrac12(p_{2,2}+s_{2,2})\quad\text{etc.}\\
&&\zeta_{2,1}&=\tfrac12(r_{2,1}-q_{2,1}),&\zeta_{2,2}&=\tfrac12(r_{2,2}-q_{2,2})\quad\text{etc.}
\end{aligned}
\]
\begin{center}and~so~on.\end{center}
\par\noindent\hspace*{\parindent}For this article, compare Articles 25---27 of Theoria interpolationis (vol. III, p. 303).]
\clearpage
\endgroup
\end{landscape}
